%------------------------------------------------------------------------------%
\section{Exercícios}

\begin{exercises*}

%------------------------------------------------------------------------------%
\begin{exercise}
  Elimine o módulo das seguintes expressões:
  \begin{alphalist}
    \item $|8| + |-3|$
    \item $|\sqrt{8}-4|$
    \item $|5x+1|$
    \item $|x^2-9|$
    \item $5-|x|$
  \end{alphalist}

  \begin{answer}
    \begin{alphalist}
      \item $11$
      \item $4-\sqrt{8}$
      \item $|5x+1|=\left\{\begin{array}{ll} 5x+1,
            \text{ se } x \geq -\frac{1}{5}\\-5x-1,
            \text{ se } x \leq -\frac{1}{5}\end{array}\right.$
      \item $|x^2-9|=\left\{\begin{array}{ll} x^2-9,
            \text{ se } x \leq -3
            \text{ ou } x \geq 3\\-(x^2-9),
            \text{ se } -3<x<3\end{array}\right.$
      \item $5-|x|=\left\{\begin{array}{ll} 5-x
            \text{ se } x \geq 0\\5+x
            \text{ se } x < 0\end{array}.\right.$
    \end{alphalist}
  \end{answer}
\end{exercise}

%------------------------------------------------------------------------------%
\begin{exercise}
  Resolva as seguintes equações modulares
  \begin{alphalist}
    \item $| 5-4x|=1$
    \item $|\frac{7-4x}{3}|=5+x$
    \item $|x-7|-|3x-1|=-x+3$
    \item $|4x^2+1|=10$
    \item $x^2+3x+1=|3x+8|$
    \item $|x|^2-8|x|+7=0$
  \end{alphalist}

  \begin{answer}
    \begin{alphalist}
      \item $S=\{ 1, \frac{3}{2}\}$
      \item $S=\{x=-\frac{8}{7}, x=22\}$
      \item $S=\{-1,\frac{5}{3}\}$
      \item $x=\pm \frac{3}{2}$
      \item $S=\{-3, \pm\sqrt{7}\}$
      \item $S=\{\pm 1, \pm 7\}$
    \end{alphalist}
  \end{answer}
\end{exercise}

%------------------------------------------------------------------------------%
\begin{exercise}
  Resolva as seguintes Inequações modulares
  \begin{alphalist}
    \item $| 5x-8| \leq 2$
    \item $|2x-3| \geq 7$
    \item $|x^2-9|+x \leq 3$
    \item $|x-2|+ |x+8| \geq 2x+12$
  \end{alphalist}

  \begin{answer}
    \begin{alphalist}
      \item $\frac{6}{5} \leq x \leq 2$
      \item $S= \{ x \in \R| x \leq -2
            \text{ ou } x \geq 5\}$
      \item $S= \{ x \in \R| -4 \leq x \leq -2
            \text{ ou } x=3\}$
      \item $x \leq -1$
    \end{alphalist}
  \end{answer}
\end{exercise}

%------------------------------------------------------------------------------%
\begin{exercise}
  Considere a função $f: \R \to \R$, definida por $f(x)= | 2x-5|$.
  \begin{alphalist}
    \item Encontre os zeros de $f$, caso existam.
    \item Verifique se $f$ é par, ímpar ou nem par nem ímpar.
    \item $f$ é injetora?
    \item Faça o esboço do gráfico de $f$.
    \item $f$ possui um valor máximo? $f$ possui valor mínimo?
  \end{alphalist}

  \begin{answer}
    \begin{alphalist}
      \item $x=\frac{5}{2}$
      \item Nem par, nem ímpar
      \item Não é injetora, pois, por exemplo, $f(2)=f(3)=1$
      \item \IncludeGraphicCentred{gabarito_aula3-4d}[0.9]
      \item O valor mínimo é 0 e não possui valor máximo
    \end{alphalist}
  \end{answer}
\end{exercise}

%------------------------------------------------------------------------------%
\begin{exercise}
  Faça o esboço do gráfico da seguintes funções:
  \begin{alphalist}
    \item $f(x)=|x^2-4x+3|$
    \item $f(x)= |x|^2-4|x|+3$
    \item $f(x)=\left\{\begin{array}{cc} 4x-3 \text{ se } x \geq
          0\\
          x^2-3x+2 \text{ se } x <0
          \end{array}\right.$
    \item $ f(x)=\left\{\begin{array}{cc} 1+x^2 \text{ se } x \leq
          0\\
          2-x \text{ se } 0 < x \leq 2 \\
          (x-2)^2  \,\,se \,\,  x> 2
          \end{array}\right.$
    \item $f(x)=|x-5|+1$
    \item $f(x)= |x-1|+ |2x-4|$
    \item $f(x)=\frac{|x|}{x}$
    \item $ f(x)=\left\{\begin{array}{cc} 1 \text{ se } x \leq
          1\\
          |x-1| \text{ se } x>1
          \end{array}\right.$
  \end{alphalist}

  \begin{answer}
    \begin{alphalist}
      \item \IncludeGraphicCentred{gabarito_aula3-5a}[0.9]
      \item \IncludeGraphicCentred{gabarito_aula3-5b}[0.9]
      \item \IncludeGraphicCentred{gabarito_aula3-5c}[0.9]
      \item \IncludeGraphicCentred{gabarito_aula3-5d}[0.9]
      \item \IncludeGraphicCentred{gabarito_aula3-5e}[0.9]
      \item \IncludeGraphicCentred{gabarito_aula3-5f}[0.9]
      \item \IncludeGraphicCentred{gabarito_aula3-5g}[0.9]
      \item \IncludeGraphicCentred{gabarito_aula3-5h}[0.9]
    \end{alphalist}
  \end{answer}
\end{exercise}

%------------------------------------------------------------------------------%
\begin{exercise}
  Dada a função $f(x)=\left\{\begin{array}{cc} 1+x^2 \text{ se } x \leq
      0\\
      2-x \text{ se } 0 < x \leq 2 \\
      (x-2^2)^2  \,\,se \,\,  x> 2
      \end{array}\right.$
  Calcule o valor da expressão $f(-2)+f(1)-f(4)$.

  \begin{answer}
    6
  \end{answer}
\end{exercise}

%------------------------------------------------------------------------------%
\begin{exercise}
  Na cidade de Pindaíba, a conta de água e esgoto tem duas faixas de preço.
  Quem consome até $10 m^3$ por mês pega um valor fixo de $R\$ 40,00$.
  Para cada metro cúbico adicional no mês, o consumidor paga $R\$6,50$.
  \begin{alphalist}
    \item Escreva a função $c(v)$ que fornece o custo mensal de água e esgoto
          em Pindaíba, com relação ao volume de água consumido (em $m^3$).
    \item Determine o valor da conta de uma casa no qual o consumo mensal
          é de $16m^3$.
    \item Trace o gráfico de $c(v)$ para $v$ entre 0 e 20 $m^3$.
  \end{alphalist}

  \begin{answer}
    \begin{alphalist}
      \item $c(v)=\left\{\begin{array}{ll} 40,
            \text{ se } v \leq 10\\40+6,5(v-10),
            \text{ se } v \geq 10\end{array}\right.$
      \item $79$
      \item \IncludeGraphicCentred{gabarito_aula3-7c}[0.9]
    \end{alphalist}
  \end{answer}
\end{exercise}

%------------------------------------------------------------------------------%
\begin{exercise}
  Duas locadoras de automóveis oferecem planos
  diferentes para a diária de um veículo econômico. A
  locadora Saturno cobra uma taxa fixa de $R\$ 30,00$, além
  de$ R\$ 0,40$ por quilômetro rodado. Já a locadora
  Mercúrio tem um plano mais elaborado: ela cobra uma
  taxa fixa de $R\$ 90,00$ com uma franquia de $200 km$, ou
  seja, o cliente pode percorrer $200 km$ sem custos adicionais.
  Entretanto, para cada km rodado além dos $200 km$
  incluídos na franquia, o cliente deve pagar $R\$ 0,60$.
  \begin{alphalist}
    \item Determine  função que descreve o custo diário de locação
          (em reais) de um automóvel em ambas locadoras.
    \item Para cada locadora, esboce o gráfico da
          função que descreve o custo diário de locação em
          termos da distância percorrida no dia.
    \item Determine para quais intervalos cada locadora tem o
          plano mais barato.
  \end{alphalist}

  \begin{answer}
    \begin{alphalist}
      \item Saturno: $f(x)=30+0,4x$,
            Mercúrio: $g(x)=\left\{\begin{array}{ll} 90,
                      \text{ se } x \leq 200 \\ 90+0,6(x-200),
                      \text{ se } x > 200\end{array}\right.$
      \item \IncludeGraphicCentred{gabarito_aula3-8b}[0.9]
      \item Saturno: $[0,150]\cup[300,\infty)$,
            Mercúrio: $[150,300]$
    \end{alphalist}
  \end{answer}
\end{exercise}

\end{exercises*}

%------------------------------------------------------------------------------%